\begin{center}
\begin{tikzpicture}[>=stealth]
 \draw (-0.5,0) -- (10,0);
 \draw[thick, domain=0:9, samples=50] plot (\x, {2*(1-(\x/9)^2)});
 \fill (0,2) circle (0.1);
 \draw[->, very thick, red] (0,2) -- (1.4,2) node[above] {$\vec v_0$};
 \draw[|<->|] (-0.4,0) -- (-0.4,2) node[midway, left] {$h$};
 \draw[|<->|] (0,-0.4) -- (9,-0.4) node[midway, below] {$w$};
\end{tikzpicture}
\end{center}

Vertikal fällt die Kugel wie beim freien Fall; ohne Luftwiderstand dauert der Flug
\begin{align}
 h &= \frac{1}{2}g\,t^2 \quad\Rightarrow\quad t = \tpF = \tp
\end{align}
Horizontal fliegt sie in dieser Zeit gleichförmig mit $v_0$:
\begin{align}
 w &= v_0\,t = \xwF \\
   &= \vp\cdot\sqrt{\frac{2\cdot\hp}{\gE}} \\
   &= \xw \approx \result{\xwP{0}{3}}
\end{align}
